\documentclass[11pt,a4paper]{article}
\usepackage{konspekt}
\renewcommand{\kheadleft}{Информатика}
\hypersetup{pdftitle={Конспект по информатике},pdfauthor={Конспект теории}}
\usetikzlibrary{automata}
\usepackage{forest}
\usepackage{listings}

\tikzset{
  fa/.style={>={Stealth[length=2.4mm,width=1.7mm]},node distance=2.4cm,auto,
    every state/.style={draw=kviolet,line width=1pt,fill=kviolet!7,minimum size=9mm,inner sep=1pt,font=\normalsize},
    accepting/.style={double=white,double distance=1.2pt,draw=korange,fill=korange!10},
    initial text={},every edge/.style={draw,line width=0.8pt},
    every loop/.style={min distance=7mm,looseness=8}},
  lbl/.style={font=\small,inner sep=1.5pt},
}
\lstdefinestyle{kcode}{basicstyle=\ttfamily\small,keywordstyle=\color{kviolet}\bfseries,
  commentstyle=\color{korange}\itshape,stringstyle=\color{korange},
  frame=leftline,framerule=2.4pt,rulecolor=\color{kviolet},backgroundcolor=\color{kviolet!4},
  xleftmargin=10pt,framexleftmargin=6pt,aboveskip=8pt,belowskip=8pt,
  columns=fullflexible,keepspaces=true,showstringspaces=false,breaklines=true,
  language=C++,morekeywords={Node,nullptr},extendedchars=true}
\lstset{style=kcode}
\lstMakeShortInline[basicstyle=\ttfamily\small\color{kviolet!70!black},breaklines=false]@
\newcommand{\tok}[1]{\tikz[baseline=(t.base)]\node[draw=kviolet,line width=0.6pt,rounded corners=1.5pt,fill=kviolet!6,inner sep=2.5pt,font=\ttfamily\small](t){#1};}
\newcommand{\rul}{\rightarrow}

\begin{document}

\makecover{Информатика}
{\begin{tikzpicture}[fa,node distance=2.6cm,scale=1.1,transform shape]
  \node[state,initial] (q0) {$q_0$};
  \node[state,right=of q0] (q1) {$q_1$};
  \node[state,accepting,right=of q1] (q2) {$q_2$};
  \node[state,below=2.2cm of q1] (q3) {$q_3$};
  \path[->] (q0) edge[kviolet] node[lbl] {$a$} (q1)
            (q1) edge[kviolet] node[lbl] {$b$} (q2)
            (q0) edge[loop above,kviolet] node[lbl] {$b$} ()
            (q2) edge[loop above,korange] node[lbl] {$a,b$} ()
            (q1) edge[bend left=20,kviolet] node[lbl] {$a$} (q3)
            (q3) edge[bend left=20,kviolet] node[lbl] {$b$} (q1)
            (q3) edge[bend right=25,korange] node[lbl,swap] {$a$} (q2);
  \begin{scope}[shift={(0.3,-4.6)},font=\ttfamily\small]
    \node[anchor=west,kviolet] at (0,0) {S ::= E \$};
    \node[anchor=west] at (0,-0.5) {E ::= T \{ + T \}};
    \node[anchor=west] at (0,-1.0) {T ::= F \{ * F \}};
    \node[anchor=west,korange] at (0,-1.5) {F ::= num | ( E )};
  \end{scope}
  \begin{scope}[shift={(6.2,-4.3)},level distance=9mm,
      every node/.style={circle,draw=kviolet,fill=white,line width=0.9pt,minimum size=6.5mm,inner sep=0pt,font=\small},
      level 1/.style={sibling distance=18mm},level 2/.style={sibling distance=10mm},
      edge from parent/.style={draw,line width=0.8pt}]
    \node[draw=korange,fill=korange!12] {$+$}
      child {node {$5$}}
      child {node[draw=korange,fill=korange!12] {$*$}
        child {node {$7$}}
        child {node {$x$}}};
  \end{scope}
\end{tikzpicture}}

\maketoc

\section{Цепочки и языки}

\textbf{Первичное понятие:} символ.

\begin{opr}
\term{Алфавит} — конечный набор символов. Обычно обозначается $\Sigma\ (V)$.
\end{opr}
\begin{opr}
\term{Цепочка} символов в алфавите $\Sigma$ — любая конечная последовательность в алфавите $\Sigma$.
\end{opr}
\begin{opr}
\term{Пустая цепочка} — цепочка, не содержащая символов. Обозначается $\varepsilon$.
\end{opr}

\textbf{Обозначения:} буквы — $a,\ b,\ c,\ 0,\ 1,\ 2$; цепочки — $\alpha,\ \beta,\ \gamma,\ \xi$.

\begin{opr}
\term{Конкатенация} цепочек $\alpha$ и $\beta$ — это $\alpha\beta$: цепочка, в которой сначала записаны символы $\alpha$, потом $\beta$.
\end{opr}
\begin{opr}
\term{Реверс} цепочки $\alpha$: $\alpha^R=a_n\cdots a_1$, если $\alpha=a_1\dots a_n$.
\end{opr}
\begin{opr}
\term{Степень} цепочки $\alpha$: $\alpha^n=\underbrace{\alpha\alpha\dots\alpha}_{n}$.
\end{opr}
\begin{opr}
$|\alpha|$ — \term{длина} цепочки (количество символов); $|\alpha|_a$ — количество символов $a$ в $\alpha$.
\end{opr}
\begin{opr}
$\Sigma^*$ — множество всех цепочек, состоящих из символов $\Sigma$. Здесь $^*$ — \term{звёздочка Клини}.
Например, при $\Sigma=\{0,1\}$:
\[
\Sigma^*=\{\varepsilon,\ 0,\ 1,\ 00,\ 01,\ 10,\ 11,\ 000,\dots\}
\]
\end{opr}
\begin{opr}
$\Sigma^+$ — множество всех \textbf{непустых} цепочек из символов $\Sigma$. Здесь $^+$ — \term{плюс Клини}, и $\Sigma^+=\Sigma^*\setminus\{\varepsilon\}$.
\end{opr}
\begin{opr}
\term{Язык} $L$ над алфавитом $\Sigma$: $L\subseteq\Sigma^*$.
\end{opr}
\begin{opr}
\term{Декартово произведение}: $A\times B=\{(a,b)\mid a\in A,\ b\in B\}$.
\end{opr}

\section{Грамматики}

\subsection{Определение грамматики и пример вывода}

\begin{opr}
\term{Грамматика} $G$ — это четвёрка объектов $G=\langle T,N,P,S\rangle$, где
\begin{itemize}
\item $T$ — множество \term{терминалов} (\emph{символы алфавита});
\item $N$ — множество \term{нетерминалов} (\emph{промежуточные символы}), причём $T\cap N=\varnothing$;
\item $P$ — множество \term{правил вывода}: $\alpha\rul\beta\in P$, где $\alpha\in(T\cup N)^+$ содержит нетерминалы, а $\beta\in(T\cup N)^*$;
\item $S$ — \term{стартовый символ}, $S\in N$.
\end{itemize}
\end{opr}

\textbf{Обозначения.} Нетерминалы: $A,\ B,\ C$.

\textbf{Пример.} $G:\ T=\{0,1\},\ N=\{S,A\}$,
\[
P:\ \begin{cases}S\rul0A1\\ 0A\rul00A1\\ A\rul\varepsilon\end{cases}
\]
Цепочка вывода:
\[
S\Rightarrow0A1\Rightarrow\overbrace{0}^{\xi_1}\overbrace{0A1}^{\gamma}\overbrace{1}^{\xi_2}\Rightarrow00\overbrace{0A1}^{\delta}11\Rightarrow\dots
\]
что даёт последовательно $01$, $0011$, $0^31^3$. Значит, $L=\{0^n1^n\mid n\ge1\}$ — язык, порождаемый грамматикой $G$.
Эквивалентная грамматика: $P':\ S\rul0S1\mid01$.

\subsection{Выводимость и сентенциальные формы}

\begin{opr}
$\beta\in(T\cup N)^*$ \term{непосредственно выводима} из $\alpha\in(T\cup N)^+$ в грамматике $G$ (обозначение $\alpha\Rightarrow_G\beta$),
если $\alpha=\xi_1\gamma\xi_2$, $\beta=\xi_1\delta\xi_2$ и правило $\gamma\rul\delta\in P$.
\end{opr}
\begin{opr}
$\beta\in(T\cup N)^*$ \term{выводима} из $\alpha\in(T\cup N)^+$ в $G$ (обозначение $\alpha\Rightarrow^*_G\beta$), если существуют
$\gamma_1,\dots,\gamma_n$ такие, что $\alpha\Rightarrow\gamma_1\Rightarrow\dots\Rightarrow\gamma_n\Rightarrow\beta$.
\end{opr}
\begin{opr}
Язык, \term{порождаемый грамматикой} $G$: $L(G)=\{\alpha\in T^*\mid S\Rightarrow^*_G\alpha\}$.
\end{opr}
\begin{opr}
\term{Сентенциальная форма} — цепочка $\alpha\in(T\cup N)^*$, выводимая из $S$.
\end{opr}

\subsection{Эквивалентность грамматик}

\begin{opr}
$G_1$ и $G_2$ называются \term{эквивалентными}, если $L(G_1)=L(G_2)$.
\end{opr}
\begin{opr}
$G_1$ и $G_2$ называются \term{почти эквивалентными}, если $L(G_1)\cup\{\varepsilon\}=L(G_2)\cup\{\varepsilon\}$.
\end{opr}

\section{Классификация Хомского}

\subsection{Типы грамматик}

\begin{itemize}
\item \textbf{Тип 0.} Все возможные.
\item \textbf{Тип 1.}
  \begin{itemize}
  \item \term{Неукорачивающие.} Правила вывода не укорачивают слово: $\alpha\rul\beta$, $|\alpha|\le|\beta|$.
  \item \term{Контекстно-зависимые.} $\alpha A\beta\rul\alpha\gamma\beta$, $|\gamma|>0$.
  \end{itemize}
\item \textbf{Тип 2.} \term{Контекстно-свободные} (КС). Все правила вывода имеют вид $A\rul\gamma$.
\item \textbf{Тип 3.}
  \begin{itemize}
  \item \term{Праволинейная}: $A\rul aB$ и $A\rul a$, где $A,B\in N$, $a\in T$.
  \item \term{Леволинейная}: $A\rul Ba$ и $A\rul a$.
  \end{itemize}
\end{itemize}

\begin{zam}
В неукорачивающих грамматиках может быть правило $S\rul\varepsilon$, если $S$ не встречается в правых частях правил.
\end{zam}

\subsection{Тип языка и обозначения}

\begin{opr}
\term{Типом языка} называется максимальный тип грамматики, задающей язык.
\end{opr}

\sidefig[0.42]{
\begin{itemize}
\item Тип 0 — формальные языки, $FL$;
\item Тип 1 — контекстно-зависимые языки, $CDL$;
\item Тип 2 — контекстно-свободные языки, $CFL$;
\item Тип 3 — регулярные языки, $REG$.
\end{itemize}
}{
\begin{tikzpicture}[scale=0.9]
  \fill[kviolet!6] (0,0) ellipse (3.6 and 1.9);   \draw[kviolet,line width=1pt] (0,0) ellipse (3.6 and 1.9);
  \fill[kviolet!13] (0.5,-0.2) ellipse (2.8 and 1.4); \draw[kviolet,line width=1pt] (0.5,-0.2) ellipse (2.8 and 1.4);
  \fill[korange!14] (0.9,-0.35) ellipse (2.0 and 0.95); \draw[korange,line width=1pt] (0.9,-0.35) ellipse (2.0 and 0.95);
  \fill[korange!32] (1.2,-0.45) ellipse (1.1 and 0.45); \draw[korange,line width=1pt] (1.2,-0.45) ellipse (1.1 and 0.45);
  \node[font=\footnotesize] at (-2.95,0) {тип 0};
  \node[font=\footnotesize] at (-1.7,-0.2) {тип 1};
  \node[font=\footnotesize] at (-0.5,-0.35) {тип 2};
  \node[font=\footnotesize] at (1.2,-0.45) {тип 3};
\end{tikzpicture}}

\subsection{Как определить тип грамматики}
\begin{enumerate}
\item Только один нетерминал слева — тип 2 или 3, иначе тип 0 или 1.
\item Все правила односторонние (нетерминал с одной стороны) — тип 3, иначе тип 2.
\end{enumerate}

\begin{zam}[Почему $a^nb^n$ контекстно-свободный, но не регулярный]
Нам нужно уметь считать, сколько $a$ мы поставили. Если ставить сначала $a$, потом $b$ по $n$ раз, число нетерминалов
не будет конечным — подходит только изменение строки из середины.
\end{zam}

\section{Структура вывода в КС-грамматиках}

\subsection{Левый и правый вывод}

\begin{opr}
\term{Левый} и \term{правый вывод} — вывод, в котором на каждом шаге раскрывается соответственно самый левый или самый правый нетерминал.
\end{opr}

Для грамматики $S\rul SS\mid a\mid b$ построим вывод слова $w=aaba$.

Левый вывод (раскрываем самый левый нетерминал):
\[
\underline{S}\rul\underline{S}S\rul\underline{S}SS\rul\underline{S}SSS\rul a\underline{S}SS\rul aa\underline{S}S\rul aab\underline{S}\rul aaba
\]
Правый вывод (раскрываем самый правый нетерминал):
\[
\underline{S}\rul S\underline{S}\rul SS\underline{S}\rul SSS\underline{S}\rul SS\underline{S}a\rul S\underline{S}ba\rul\underline{S}aba\rul aaba
\]

\subsection{Дерево вывода. Неоднозначность}

\begin{opr}
\term{Дерево вывода} — дерево, узлами которого являются нетерминалы, а листьями — терминалы. Дерево вывода не единственно.
\end{opr}

\begin{figure}[H]
\centering
\begin{forest}
  for tree={circle,draw=kviolet,fill=kviolet!7,line width=0.9pt,minimum size=7mm,inner sep=0pt,
    l sep=5mm,s sep=6mm,edge={line width=0.8pt}},
  [$S$ [$S$ [$S$ [$S$ [$a$,draw=korange,fill=korange!12]] [$S$ [$a$,draw=korange,fill=korange!12]]] [$S$ [$b$,draw=korange,fill=korange!12]]]
       [$S$ [$a$,draw=korange,fill=korange!12]]]
\end{forest}
\caption{Дерево вывода слова $aaba$ в грамматике $S\rul SS\mid a\mid b$}
\end{figure}

\begin{opr}
Грамматика называется \term{неоднозначной}, если у неё существует несколько деревьев вывода для некоторой цепочки.
В ином случае грамматика называется \term{однозначной}.
\end{opr}

\subsection{Недостижимые и бесплодные символы}

\begin{opr}
\term{Недостижимые символы} — нетерминалы, которые не выводятся из $S$.
\end{opr}
\begin{opr}
\term{Бесплодные символы} — нетерминалы, из которых не выводятся терминальные цепочки.
\end{opr}

\textbf{Пример.}
\[
\begin{array}{l}
S\rul aAB\mid E\\ A\rul dDA\\ B\rul bE\mid f\\ C\rul cAB\mid dSD\mid a\\ D\rul eA\\ E\rul fA\mid g
\end{array}
\]
Недостижимые: $C$. Бесплодные: $A,\ D$.

\section{Преобразование грамматик}

\subsection{Очистка грамматики}

\begin{kvb}{Поиск}
\begin{enumerate}
\item \textbf{Недостижимые:}
  \begin{enumerate}
  \item строим $K_0=\{S\}$;
  \item строим $K_i=K_{i-1}\cup\{B\mid\exists A\rul\xi_1B\xi_2,\ A\in K_{i-1}\}$;
  \item конец, когда $K_i=K_{i-1}$.
  \end{enumerate}
\item \textbf{Бесплодные:}
  \begin{enumerate}
  \item строим $K_0=\Sigma$ (то есть $T$);
  \item строим $K_i=K_{i-1}\cup\{A\mid\exists A\rul\gamma,\ \gamma\in K_{i-1}^*\}$;
  \item конец, когда $K_i=K_{i-1}$.
  \end{enumerate}
\end{enumerate}
\end{kvb}

\begin{kob}{Удаление}
\begin{enumerate}
\item Удаляем бесплодные.
\item Удаляем недостижимые.
\end{enumerate}
\end{kob}

\textbf{Пример а).}
\[
\begin{array}{l}
S\rul aABS\mid bCACd\\ A\rul bAB\mid cSA\mid cCC\\ B\rul bAB\mid cSB\\ C\rul cS\mid c
\end{array}
\]
Плодовитые:
\[
\begin{array}{l}
K_0=\{a,b,c,d\}\\ K_1=\{a,b,c,d,C\}\\ K_2=\{a,b,c,d,C,A\}\\ K_3=\{a,b,c,d,C,A,S\}\\ K_4=K_3\Rightarrow\text{конец}
\end{array}
\]
$\Rightarrow$ бесплодные: $B$. Тогда
\[
P:\ \begin{array}{l}S\rul bCACd\\ A\rul cSA\mid cCC\\ C\rul cS\mid c\end{array}
\]
Достижимые: $K_0=\{S\}$, $K_1=\{S,C,A\}$, $K_2=K_1\Rightarrow$ конец $\Rightarrow$ недостижимых нет.

\textbf{Пример б).}
\[
\begin{array}{l}
S\rul ab\mid aA\\ A\rul bAB\mid aA\\ B\rul aB\mid b
\end{array}
\]
Плодовитые: $K_0=\{a,b\}$, $K_1=\{a,b,S,B\}$, $K_2=K_1\Rightarrow$ конец $\Rightarrow$ бесплодные: $A$. Тогда
$P:\ S\rul ab,\ B\rul aB\mid b$.

Достижимые: $K_0=\{S\}$, $K_1=K_0\Rightarrow$ конец $\Rightarrow$ недостижимые: $B$, и остаётся $P:\ S\rul ab$.

\subsection{Избавление от \texorpdfstring{$\varepsilon$}{ε}-правил}

\[
\begin{array}{l}A\rul aB\mid\varepsilon\\ C\rul DABA\end{array}
\qquad\Longrightarrow\qquad
\begin{array}{l}A\rul aB\\ C\rul DABA\mid DB\mid DAB\mid DBA\end{array}
\]

\begin{kvb}{Алгоритм}
\begin{enumerate}
\item Строим $N'$ — множество \term{«схлопывающихся»} нетерминалов: $N'=\{A\in N\mid A\rul^*\varepsilon\}$.
\item Убираем правила вида $A\rul\varepsilon$. Если $\varepsilon\in L$, то вводим новый стартовый символ $S'$ с правилом $S'\rul S\mid\varepsilon$.
\item Для всех правил вида $B\rul\alpha_0A_1\alpha_1A_2\alpha_2\dots A_n\alpha_n$, где $A_i\in N'$ и $\alpha_i\in(T\cup(N\setminus N'))^*$,
  добавляем правила со всеми возможными выбрасываниями $A_i$:
  \[
  \begin{array}{l}
  B\rul\alpha_0\alpha_1\dots\alpha_n\\ B\rul\alpha_0\alpha_1\dots A_n\alpha_n\\ \dots\\ B\rul\alpha_0A_1\alpha_1A_2\dots A_n\alpha_n
  \end{array}
  \]
\item Удалить бесплодные и недостижимые.
\end{enumerate}
\end{kvb}

\section{Компиляция на практике}

\begin{kob}{Как это используется на практике}
Для компиляторов используют регулярные (тип 3) и контекстно-свободные (тип 2) языки.
\end{kob}

\subsection{Этапы компиляции}

\begin{enumerate}
\item \textbf{Лексический анализ.} Текст $\rul$ поток токенов:
  \[
  \texttt{if(a<5)}\ \rul\
  \underset{\text{служеб. слово}}{\tok{if}}\ \tok{(}\ \underset{\text{имя}}{\tok{a}}\ \underset{\text{опер.}}{\tok{<}}\ \underset{\text{число}}{\tok{5}}\ \tok{)}
  \]
\item \textbf{Синтаксический анализ.} Поток токенов $\rul$ синтаксическое дерево (рис.~\ref{fig:if}).
\item \textbf{Семантический анализ.} Проверка на смысловую корректность.
\item \textbf{Генерация кода.}
\end{enumerate}

\begin{figure}[H]
\centering
\begin{forest}
  for tree={draw=kviolet,fill=kviolet!7,line width=0.9pt,rounded corners=2pt,font=\ttfamily,
    minimum height=6mm,minimum width=8mm,inner sep=2pt,l sep=6mm,s sep=8mm,edge={line width=0.8pt}},
  [if,draw=korange,fill=korange!12 [<\vphantom{p} [a] [5]] [тело,font=\normalfont\itshape,draw=kgray,fill=white,dashed]]
\end{forest}
\caption{Синтаксическое дерево для \texttt{if(a<5)}}\label{fig:if}
\end{figure}

\subsection{Подходы к компиляции}

\begin{enumerate}
\item \textbf{Традиционная компиляция.} $\texttt{file.cpp}\rul\texttt{file.exe}$ (можно сразу запустить). Быстрые компиляция и работа.
\item \textbf{JIT-компиляция} (JIT — Just In Time). $\texttt{file.cs}\rul\text{IL}$ (промежуточный язык), компиляция в машинные коды
  при запуске. Медленная компиляция, быстрая работа, безопасность.
\item \textbf{Интерпретация.} \texttt{file.py}. Медленная работа, кроссплатформенность, безопасность.
\end{enumerate}

\subsection{Распознающие формализмы}

Грамматика — \term{порождающий} формализм. Ему двойственны \term{распознающие}:

\begin{table}[H]
\centering
\renewcommand{\arraystretch}{1.35}
\begin{tabular}{ll}
\toprule
{\hfont\color{kviolet}Тип языка} & {\hfont\color{kviolet}Распознающий формализм}\\
\midrule
Тип 3 & Конечные автоматы\\
Тип 2 & МП-автоматы (автоматы с магазинной памятью)\\
Тип 1 & ЛОА (линейно ограниченный автомат)\\
Тип 0 & Машина Тьюринга\\
\bottomrule
\end{tabular}
\caption{Порождающие и распознающие формализмы}
\end{table}

\section{Регулярные языки}

\subsection{НКА и ДКА}

\begin{opr}
\term{Недетерминированный конечный автомат (НКА)} — пятёрка объектов $\langle Q,\Sigma,\delta,q_0,F\rangle$:
\begin{itemize}
\item $Q$ — конечное множество состояний;
\item $\Sigma$ — алфавит;
\item $\delta:\ Q\times(\Sigma\cup\{\varepsilon\})\rul2^Q$ (множество всех подмножеств $Q$) — функция перехода (переход не однозначен);
\item $q_0\in Q$ — начальное состояние;
\item $F\subseteq Q$ — множество конечных состояний.
\end{itemize}
\end{opr}

\begin{opr}
\term{Детерминированный конечный автомат (ДКА)} — пятёрка объектов $\langle Q,\Sigma,\delta,q_0,F\rangle$:
\begin{itemize}
\item $Q$ — конечное множество состояний;
\item $\Sigma$ — алфавит;
\item $\delta:\ Q\times\Sigma\rul Q\cup\{\varnothing\}$ — функция перехода (переход всегда однозначен);
\item $q_0\in Q$ — начальное состояние;
\item $F\subseteq Q$ — множество конечных состояний.
\end{itemize}
\end{opr}

\begin{figure}[H]
\centering
\begin{tikzpicture}[fa,node distance=3cm]
  \node[state,initial,label={[kviolet,font=\small]above:$S$}] (q0) {$q_0$};
  \node[state,right=of q0,label={[kviolet,font=\small]above left:$B$}] (q1) {$q_1$};
  \node[state,below=2.2cm of q0,label={[kviolet,font=\small]left:$A$}] (q2) {$q_2$};
  \node[state,right=of q2,label={[kviolet,font=\small]below:$C$}] (q3) {$q_3$};
  \node[state,accepting,right=3.2cm of q3] (q4) {$q_4$};
  \node[draw=korange,text=korange,rounded corners=2pt,font=\hfont\small,right=1.4cm of q1] {НКА};
  \path[->] (q0) edge node[lbl] {$0,1$} (q1)
            (q0) edge node[lbl,swap] {$0$} (q2)
            (q2) edge[bend left=12] node[lbl] {$1$} (q1)
            (q1) edge[bend left=12] node[lbl] {$0$} (q2)
            (q1) edge[loop above] node[lbl] {$0$} ()
            (q1) edge node[lbl] {$1$} (q3)
            (q2) edge[loop below] node[lbl] {$1$} ()
            (q2) edge node[lbl,swap] {$0$} (q3)
            (q3) edge[korange] node[lbl] {$\perp$} (q4);
  \node[font=\footnotesize,korange,below=2pt] at ($(q3)!0.5!(q4)$) {терминальный переход};
\end{tikzpicture}
\caption{Автомат к примеру 1}
\end{figure}

\begin{minipage}[t]{0.48\linewidth}
Праволинейная грамматика:
\[
\begin{array}{l}
S\rul0A\mid0B\mid1B\\ A\rul1A\mid1B\mid0C\\ B\rul0A\mid0B\mid1C\\ C\rul\perp
\end{array}
\]
\end{minipage}\hfill
\begin{minipage}[t]{0.48\linewidth}
Леволинейная грамматика:
\[
\begin{array}{l}
D\rul C\perp\\ C\rul B1\mid A0\\ B\rul B0\mid A1\mid0\mid1\\ A\rul A1\mid B0\mid0
\end{array}
\]
\end{minipage}

\begin{figure}[H]
\centering
\begin{tikzpicture}[fa,node distance=3.2cm]
  \node[state,initial] (S) {$S$};
  \node[state,accepting,right=of S] (A) {$A$};
  \node[state,right=of A] (B) {$B$};
  \path[->] (S) edge[bend left=20] node[lbl] {$a$} (A)
            (A) edge[bend left=20] node[lbl] {$b$} (S)
            (A) edge[bend left=20] node[lbl] {$a$} (B)
            (B) edge[bend left=20] node[lbl] {$b$} (A)
            (S) edge[loop below] node[lbl] {$b$} ()
            (B) edge[loop below] node[lbl] {$b$} ();
\end{tikzpicture}
\caption{Автомат к примеру 2}
\end{figure}

\begin{minipage}[t]{0.48\linewidth}
Праволинейная грамматика:
\[
\begin{array}{l}
S\rul aA\mid bS\mid a\\ A\rul bS\mid aB\\ B\rul bB\mid bA\mid b
\end{array}
\]
\end{minipage}\hfill
\begin{minipage}[t]{0.48\linewidth}
Леволинейная грамматика:
\[
\begin{array}{l}
A\rul Sa\mid Bb\mid a\\ B\rul Ba\mid Aa\\ S\rul Sb\mid Ab\mid b
\end{array}
\]
\end{minipage}

\subsection{Автоматы и линейные грамматики}

Пусть есть автомат $\langle Q,\Sigma,\gamma,q_0,F\rangle$ и грамматика $\langle T,N,S,P\rangle$. Тогда можно сопоставить:

\begin{minipage}[t]{0.48\linewidth}
\textbf{Для праволинейных:}
\begin{itemize}
\item $Q\rul N$, \quad $\Sigma\rul T$,
\item $q_0\rul S$, \quad $F\rul -$,
\item $\gamma\rul P$.
\end{itemize}
\end{minipage}\hfill
\begin{minipage}[t]{0.48\linewidth}
\textbf{Для леволинейных:}
\begin{itemize}
\item $Q\rul N$, \quad $\Sigma\rul T$,
\item $q_0\rul -$, \quad $F\rul S$,
\item $\gamma\rul P$.
\end{itemize}
\end{minipage}

\begin{figure}[H]
\centering
\begin{tikzpicture}[fa,node distance=1.9cm]
  \node[state] (A1) {$A$}; \node[state,right=of A1,dashed] (B1) {$B$};
  \node[right=1.4cm of B1] (r1) {$A\rul\alpha B$};
  \node[state,below=1.2cm of A1] (A2) {$A$}; \node[state,accepting,right=of A2] (B2) {$B$};
  \node[right=1.4cm of B2] (r2) {$A\rul\alpha$};
  \path[->] (A1) edge node[lbl] {$\alpha$} (B1) (A2) edge node[lbl] {$\alpha$} (B2);
  \node[font=\Large,kviolet] at ($(B1)!0.5!(B2)+(1.0,0)$) {$\Longleftrightarrow$};
  \node[font=\hfont\small,kviolet] at ($(A1)!0.5!(r1)+(0,1.0)$) {праволинейная};
  \begin{scope}[xshift=8.2cm]
  \node[state] (C1) {$A$}; \node[state,right=of C1] (D1) {$B$};
  \draw[->,dashed,kgray] ($(C1)+(-1.2,0)$) -- (C1);
  \node[right=1.4cm of D1] (s1) {$B\rul A\alpha$};
  \node[state,initial,below=1.2cm of C1] (C2) {$A$}; \node[state,accepting,right=of C2] (D2) {$B$};
  \node[right=1.4cm of D2] (s2) {$B\rul\alpha$};
  \path[->] (C1) edge node[lbl] {$\alpha$} (D1) (C2) edge node[lbl] {$\alpha$} (D2);
  \node[font=\Large,korange] at ($(D1)!0.5!(D2)+(1.0,0)$) {$\Longleftrightarrow$};
  \node[font=\hfont\small,korange] at ($(C1)!0.5!(s1)+(0,1.0)$) {леволинейная};
  \end{scope}
\end{tikzpicture}
\caption{Соответствие автомата и линейной грамматики (пунктирное состояние — любое)}
\end{figure}

\needspace{10\baselineskip}\textbf{Пример 3.}
\sidefig[0.55]{
\[
\begin{array}{l}
S\rul aS\mid bN\\ N\rul aN\mid aK\mid b\\ K\rul a\mid bK
\end{array}
\]
}{
\begin{tikzpicture}[fa,node distance=2.2cm]
  \node[state,initial] (S) {$S$};
  \node[state,right=of S] (N) {$N$};
  \node[state,right=of N] (K) {$K$};
  \node[state,accepting,below=1.6cm of N] (F) {$\perp$};
  \path[->] (S) edge node[lbl] {$b$} (N) (N) edge node[lbl] {$a$} (K)
            (N) edge node[lbl,swap] {$b$} (F) (K) edge node[lbl] {$a$} (F)
            (S) edge[loop above] node[lbl] {$a$} () (N) edge[loop above] node[lbl] {$a$} () (K) edge[loop above] node[lbl] {$b$} ();
\end{tikzpicture}}

\needspace{10\baselineskip}\textbf{Пример 4.}
\sidefig[0.55]{
\[
\begin{array}{l}
S\rul aN\mid kS\mid cS\\ N\rul a\mid cN\mid cA\mid k\\ A\rul b\mid bN\mid eA\mid d\mid cS
\end{array}
\]
}{
\begin{tikzpicture}[fa,node distance=2.4cm]
  \node[state,initial] (S) {$S$};
  \node[state,right=of S] (N) {$N$};
  \node[state,below=1.9cm of $(S)!0.5!(N)$] (A) {$A$};
  \node[state,accepting,right=2.6cm of A] (F) {$\perp$};
  \path[->] (S) edge node[lbl] {$a$} (N)
            (N) edge[bend left=15] node[lbl] {$c$} (A)
            (A) edge[bend left=15] node[lbl] {$b$} (N)
            (A) edge[bend left=15] node[lbl] {$c$} (S)
            (N) edge node[lbl] {$a,k$} (F)
            (A) edge node[lbl,swap] {$b,d$} (F)
            (S) edge[loop above] node[lbl] {$k,c$} ()
            (N) edge[loop above] node[lbl] {$c$} ()
            (A) edge[loop below] node[lbl] {$e$} ();
\end{tikzpicture}}

\needspace{10\baselineskip}\textbf{Пример 5} (леволинейная грамматика).
\sidefig[0.55]{
\[
\begin{array}{l}
S\rul S0\mid S1\mid P0\mid P1\\ P\rul N.\\ N\rul0\mid1\mid N0\mid N1
\end{array}
\]
}{
\begin{tikzpicture}[fa,node distance=1.5cm]
  \node[state,initial] (H) {$\perp$};
  \node[state,right=of H] (N) {$N$};
  \node[state,right=of N] (P) {$P$};
  \node[state,accepting,right=of P] (S) {$S$};
  \path[->] (H) edge node[lbl] {$0,1$} (N) (N) edge node[lbl] {$.$} (P) (P) edge node[lbl] {$0,1$} (S)
            (N) edge[loop above] node[lbl] {$0,1$} () (S) edge[loop above] node[lbl] {$0,1$} ();
\end{tikzpicture}}

\subsection{Построение ДКА}

Разберём на примере 3. Состояния ДКА — множества состояний НКА:

\begin{table}[H]
\centering
\renewcommand{\arraystretch}{1.3}
\begin{tabular}{c|c|c|c}
\toprule
{\hfont\color{kviolet}ДКА} & {\hfont\color{kviolet}НКА} & $a$ & $b$\\
\midrule
$Q_0$ & $S$ & $Q_0$ & $Q_1$\\
$Q_1$ & $N$ & $Q_2$ & $Q_3$\\
$Q_2$ & $N,K$ & $Q_4$ & $Q_5$\\
$Q_3$ & $\perp$ & — & —\\
$Q_4$ & $N,K,\perp$ & $Q_4$ & $Q_5$\\
$Q_5$ & $K,\perp$ & $Q_3$ & $Q_6$\\
$Q_6$ & $K$ & $Q_3$ & $Q_6$\\
\bottomrule
\end{tabular}
\caption{Построение ДКА по НКА из примера 3}
\end{table}

\needspace{10\baselineskip}\textbf{Пример 6.} НКА задан таблицей переходов:

\begin{minipage}[c]{0.45\linewidth}\centering
\renewcommand{\arraystretch}{1.3}
\begin{tabular}{c|c|c|c|c}
\toprule
 & $\varepsilon$ & $a$ & $b$ & $c$\\
\midrule
$\rul p$ & $\varnothing$ & $\{p\}$ & $\{q\}$ & $\{r\}$\\
$q$ & $\{p\}$ & $\{q\}$ & $\{r\}$ & $\varnothing$\\
$^*r$ & $\{q\}$ & $\{r\}$ & $\varnothing$ & $\{p\}$\\
\bottomrule
\end{tabular}
\end{minipage}\hfill
\begin{minipage}[c]{0.5\linewidth}\centering
\begin{tikzpicture}[fa,node distance=2.6cm]
  \node[state,initial] (p) {$p$};
  \node[state,right=of p] (q) {$q$};
  \node[state,accepting,below=1.8cm of $(p)!0.5!(q)$] (r) {$r$};
  \path[->] (p) edge node[lbl] {$b$} (q)
            (q) edge[bend right=35] node[lbl,swap] {$\varepsilon$} (p)
            (p) edge[bend left=12] node[lbl] {$c$} (r)
            (r) edge[bend left=25] node[lbl] {$c$} (p)
            (q) edge[bend left=12] node[lbl] {$b$} (r)
            (r) edge[bend right=25] node[lbl,swap] {$\varepsilon$} (q)
            (p) edge[loop above] node[lbl] {$a$} ()
            (q) edge[loop above] node[lbl] {$a$} ()
            (r) edge[loop below] node[lbl] {$a$} ();
\end{tikzpicture}
\end{minipage}

\medskip
\begin{minipage}[c]{0.45\linewidth}\centering
\renewcommand{\arraystretch}{1.3}
\begin{tabular}{c|c|c|c|c}
\toprule
{\hfont\color{kviolet}ДКА} & {\hfont\color{kviolet}НКА} & $a$ & $b$ & $c$\\
\midrule
$\rul Q_0$ & $p$ & $Q_0$ & $Q_1$ & $Q_2$\\
$Q_1$ & $q,p$ & $Q_1$ & $Q_2$ & $Q_2$\\
$^*Q_2$ & $q,p,r$ & $Q_2$ & $Q_2$ & $Q_2$\\
\bottomrule
\end{tabular}
\end{minipage}\hfill
\begin{minipage}[c]{0.5\linewidth}\centering
\begin{tikzpicture}[fa,node distance=2.6cm]
  \node[state,initial] (Q0) {$Q_0$};
  \node[state,right=of Q0] (Q1) {$Q_1$};
  \node[state,accepting,below=1.8cm of $(Q0)!0.5!(Q1)$] (Q2) {$Q_2$};
  \path[->] (Q0) edge node[lbl] {$b$} (Q1) (Q0) edge node[lbl,swap] {$c$} (Q2) (Q1) edge node[lbl] {$b,c$} (Q2)
            (Q0) edge[loop above] node[lbl] {$a$} () (Q1) edge[loop above] node[lbl] {$a$} () (Q2) edge[loop below] node[lbl] {$a,b,c$} ();
\end{tikzpicture}
\end{minipage}

\needspace{8\baselineskip}\textbf{Пример 7.}
\begin{figure}[H]
\centering
\begin{tikzpicture}[fa,node distance=3cm]
  \node[state,initial] (S) {$S$};
  \node[state,right=of S] (B) {$B$};
  \node[state,accepting,below=1.9cm of S] (A) {$A$};
  \node[state,below=1.9cm of B] (C) {$C$};
  \node[state,accepting,right=2.8cm of $(B)!0.5!(C)$] (D) {$D$};
  \path[->] (S) edge[bend left=15] node[lbl] {$a$} (B)
            (S) edge[bend right=15] node[lbl,swap] {$\varepsilon$} (B)
            (S) edge node[lbl,swap] {$a$} (A)
            (B) edge node[lbl,swap,pos=0.35] {$b$} (A)
            (B) edge node[lbl] {$b$} (C)
            (B) edge node[lbl] {$a$} (D)
            (A) edge node[lbl,swap] {$a$} (C)
            (C) edge node[lbl,swap] {$\varepsilon$} (D)
            (S) edge[loop above] node[lbl] {$b$} ()
            (B) edge[loop above] node[lbl] {$b$} ();
\end{tikzpicture}
\caption{НКА к примеру 7}
\end{figure}

\begin{minipage}[c]{0.45\linewidth}\centering
\renewcommand{\arraystretch}{1.3}
\begin{tabular}{c|c|c|c}
\toprule
{\hfont\color{kviolet}ДКА} & {\hfont\color{kviolet}НКА} & $a$ & $b$\\
\midrule
$\rul Q_0$ & $SB$ & $Q_1$ & $Q_2$\\
$^*Q_1$ & $ABD$ & $Q_3$ & $Q_4$\\
$^*Q_2$ & $SABCD$ & $Q_4$ & $Q_2$\\
$^*Q_3$ & $CD$ & $\varnothing$ & $\varnothing$\\
$^*Q_4$ & $ABCD$ & $Q_3$ & $Q_4$\\
\bottomrule
\end{tabular}
\end{minipage}\hfill
\begin{minipage}[c]{0.52\linewidth}\centering
\begin{tikzpicture}[fa,node distance=2.3cm]
  \node[state,initial] (Q0) {$Q_0$};
  \node[state,accepting,right=of Q0] (Q1) {$Q_1$};
  \node[state,accepting,right=of Q1] (Q3) {$Q_3$};
  \node[state,accepting,below=1.6cm of Q0] (Q2) {$Q_2$};
  \node[state,accepting,below=1.6cm of Q1] (Q4) {$Q_4$};
  \path[->] (Q0) edge node[lbl] {$a$} (Q1) (Q1) edge node[lbl] {$a$} (Q3)
            (Q0) edge node[lbl,swap] {$b$} (Q2) (Q1) edge node[lbl] {$b$} (Q4)
            (Q2) edge node[lbl,swap] {$a$} (Q4) (Q4) edge node[lbl,swap] {$a$} (Q3)
            (Q2) edge[loop below] node[lbl] {$b$} () (Q4) edge[loop below] node[lbl] {$b$} ();
\end{tikzpicture}
\end{minipage}

\medskip
\begin{minipage}[t]{0.48\linewidth}
Праволинейная грамматика (убрали $Q_3$, т.\,к. бесплодное):
\[
\begin{array}{l}
Q_0\rul aQ_1\mid bQ_2\mid a\mid b\\ Q_1\rul bQ_4\mid a\mid b\\ Q_2\rul aQ_4\mid bQ_2\mid a\mid b\\ Q_4\rul bQ_4\mid a\mid b
\end{array}
\]
\end{minipage}\hfill
\begin{minipage}[t]{0.48\linewidth}
Леволинейная грамматика:
\[
\begin{array}{l}
\perp\rul Q_1\mid Q_2\mid Q_3\mid Q_4\\ Q_3\rul Q_1a\mid Q_4a\\ Q_4\rul Q_4b\mid Q_2a\mid Q_1b\\ Q_2\rul Q_2b\mid b\\ Q_1\rul a
\end{array}
\]
\end{minipage}

\subsection{Операции над языками и регулярные выражения}

\begin{opr}
\term{Конкатенация} языков: $AB=\{\alpha\beta\mid\alpha\in A,\ \beta\in B\}$.
\end{opr}
\textbf{Пример.} $A=\{a,bb\}$, $B=\{ba,cbc\}$ $\Rightarrow$ $AB=\{aba,\ acbc,\ bbba,\ bbcbc\}$.

\begin{opr}
\term{Степень} языка: $L^n=\underbrace{LL\dots L}_{n}$, $L^0=\{\varepsilon\}$.
\end{opr}
\textbf{Пример.} $A=\{a,bb\}$: $A^3=\{aaa,\ aabb,\ abba,\ abbbb,\ bbaa,\ bbabb,\ bbbba,\ bbbbbb\}$.

\begin{opr}
\term{Замыкание Клини}: $L^*=\bigcup\limits_{n=0}^{\infty}L^n$.
\end{opr}
\textbf{Пример.} $A=\{a,bb\}$: $A^*=\{\varepsilon,\ a,\ bb,\ aa,\ abb,\ bba,\ bbbb,\dots\}$.

\begin{opr}
\term{Плюс Клини}: $L^+=\bigcup\limits_{n=1}^{\infty}L^n=LL^*$.
\end{opr}

\begin{opr}
\term{Регулярное выражение} состоит из: $\varepsilon,\ \Sigma,\ (,\ ),\ ^*,\ |,\ \cdot,\ \varnothing$.
\begin{enumerate}
\item $\varnothing$, $\varepsilon$, $a$ ($a\in\Sigma$) — регулярные выражения.
\item Если $\alpha$ и $\beta$ — регулярные, то $\alpha|\beta$, $\alpha\cdot\beta$, $\alpha^*$, $(\alpha)$ — регулярные.
\item Других регулярных выражений нет.
\end{enumerate}
\end{opr}

\begin{opr}
Язык, задаваемый \term{регулярным выражением}:
\begin{enumerate}
\item $L(\varnothing)=\varnothing$, \quad $L(\varepsilon)=\{\varepsilon\}$, \quad $L(a)=\{a\}$;
\item $L(\alpha\cdot\beta)=L(\alpha)L(\beta)$, \quad $L(\alpha^*)=L(\alpha)^*$, \quad $L((\alpha))=L(\alpha)$, \quad $L(\alpha|\beta)=L(\alpha)\cup L(\beta)$.
\end{enumerate}
\end{opr}

\textbf{Пример.} Язык всех слов над $\{a,b\}$, заканчивающихся на $abb$. Регулярное выражение: $(a|b)^*\cdot a\cdot b\cdot b$.

\subsection{Регулярные выражения на практике (PCRE)}

Регулярное выражение — шаблон для поиска, проверки и замены текста. В библиотеках синтаксис богаче, чем в определении выше:
есть обратные ссылки, которые выходят за рамки регулярных языков.

\textbf{Диалекты:} POSIX (BRE, ERE; @grep@, @sed@), PCRE (Perl-Compatible Regular Expressions, самый богатый синтаксис).
В C++ @std::regex@ использует ECMAScript, близкий к PCRE (но без lookbehind).

\textbf{Функции в C++} (@<regex>@):
\begin{itemize}
\item @regex_match@ — проверка \textbf{полного} совпадения строки;
\item @regex_search@ — поиск \textbf{первого} вхождения;
\item @regex_replace@ — замена совпадающих частей.
\end{itemize}
Результат лежит в @std::smatch@ / @std::cmatch@: @res[0]@ — всё совпадение, @res[1]@, @res[2]@, \ldots{} — группы.

\begin{zam}
В строке C++ обратный слеш удваивается, проще писать raw-строку: @"\\d+"@ $\Rightarrow$ @R"(\d+)"@.
\end{zam}

\begin{opr}
\term{Метасимволы} — символы со специальным смыслом:
\begin{center}@. ^ $ * + ? ( ) [ ] { } | \@\end{center}
Чтобы использовать метасимвол как обычный, его \term{экранируют}: @\.@, @\*@, @\\@.
\end{opr}

\begin{opr}
\term{Одиночные метасимволы}:
\begin{itemize}
\item @.@ — любой символ, кроме @\n@;
\item @^@ — начало строки, @$@ — конец строки;
\item @\b@ — граница слова, @\B@ — не граница слова.
\end{itemize}
\end{opr}

\begin{zam}
@\b@ и @\w@ работают только для латиницы.
\end{zam}

\begin{opr}
\term{Символьные классы} (заглавная буква — инверсия):
\begin{itemize}
\item @\s@ (@\S@) — пробельные символы;
\item @\d@ (@\D@) — цифры 0–9;
\item @\w@ (@\W@) — буквы, цифры, знак подчёркивания.
\end{itemize}
\end{opr}

\begin{opr}
\term{Перечисление} @[...]@ — один любой символ из набора, внутри можно использовать диапазоны: @[a-fA-F0-5]@.
\term{Инвертированное} перечисление: @[^abc]@ — любой символ, кроме перечисленных. Внутри @[]@ метасимволы означают сами себя:
@[.]@ — просто точка.
\end{opr}

\begin{opr}
\term{Квантификаторы} — сколько раз повторять предыдущий символ или группу:
\begin{itemize}
\item @*@ — ноль или более повторений; @+@ — одно или более; @?@ — ноль или одно;
\item @{n}@ — ровно $n$ раз; @{n,m}@ — от $n$ до $m$ раз; @{n,}@ — от $n$ до $\infty$ раз.
\end{itemize}
\end{opr}

\begin{zam}
Формы @{,m}@ нет: @a{,5}@ — это просто текст, для «от 0 до 5» пишем @a{0,5}@.
\end{zam}

\begin{opr}
\term{Типы квантификации}:
\begin{itemize}
\item \textbf{жадная} (@+@, @*@, @?@, @{}@) — берёт самую длинную подстроку;
\item \textbf{ленивая} (после квантификатора ставим @?@: @+?@, @*?@, @??@) — берёт самую короткую.
\end{itemize}
\end{opr}

Например, на строке @<b>x</b>@ шаблон @<.+>@ захватит всю строку, а @<.+?>@ — только @<b>@.

\subsection{Группировки}

\begin{opr}
\term{Группировки} ставятся в круглые скобки:
\begin{itemize}
\item @(...)@ — \textbf{группа с захватом}: запоминает подстроку (достаём по индексу);
\item @(?:...)@ — \textbf{группа без захвата}: ничего не запоминает, просто делит шаблон на логические блоки;
\item @|@ — «или», самый низкий приоритет, поэтому обычно берётся в группу.
\end{itemize}
Пример: @(?:red|green) light@ (без скобок @red|green light@ значит «red» или «green light»).
Группы нумеруются по \textbf{открывающей} скобке слева направо, @res[0]@ — всё совпадение.
\end{opr}

\begin{opr}
\term{Просмотр вперёд и назад} (lookaround) — проверка контекста, который не входит в результат:
\begin{itemize}
\item @(?=string)@ — дальше должно быть @string@;
\item @(?!string)@ — дальше @string@ быть не должно;
\item @(?<=string)@, @(?<!string)@ — то же, но назад (в @std::regex@ не поддерживается).
\end{itemize}
\end{opr}

\textbf{Пример.} Шаблон @best(?= match)@ на строке «best match» даёт @res[0]@ $=$ «best»: «match» проверили, но в результат не взяли.

\begin{opr}
\term{Обратная ссылка} (backreference) @\num@ — ссылка на группу захвата с номером @num@: ровно тот же текст, который она захватила.
\end{opr}

\textbf{Пример.} Шаблон @(red|green) apple is \1@ подойдёт к «red apple is red» и «green apple is green», но не к «green apple is red».

\begin{table}[H]
\centering
\renewcommand{\arraystretch}{1.4}
\begin{tabular}{p{0.24\linewidth}p{0.32\linewidth}p{0.32\linewidth}}
\toprule
{\hfont\color{kviolet}Задача} & {\hfont\color{kviolet}Шаблон} & {\hfont\color{kviolet}Замечание}\\
\midrule
Файл-картинка & @.*\.(jpg|png|gif)@ & точка экранирована\\
Дата ДД-ММ-ГГГГ & @\d{2}-\d{2}-\d{4}@ & то же, что @[0-9]{2}-@\ldots\\
HEX-цвет & @#[A-Fa-f0-9]{6}@ & не забыть строчные @a-f@\\
\bottomrule
\end{tabular}
\caption{Практические шаблоны}
\end{table}

\begin{kob}{Частые ошибки}
\begin{itemize}
\item путать @match@ (вся строка) и @search@ (подстрока);
\item забывать удваивать @\@ в обычных строках C++;
\item ждать от @{,n}@ смысла «до $n$»;
\item не брать @|@ в скобки: @red|green light@ $\ne$ @(?:red|green) light@.
\end{itemize}
\end{kob}

\subsection{Лексический анализ на практике}

Токены образуют регулярный язык $\Rightarrow$ лексический анализатор строится как конечный автомат (обычно через рекурсивные функции).

\begin{kvb}{Требования к лексическому анализатору}
\begin{enumerate}
\item время работы $O(n)$, где $n$ — длина входного текста;
\item реализация через конечный автомат;
\item поиск ключевых слов по бору, а не перебором списка.
\end{enumerate}
\end{kvb}

\textbf{Типы токенов:}
\begin{enumerate}
\item ключевые слова (key words);
\item операторы (operator);
\item литералы (literal);
\item идентификаторы, т.\,е. имена переменных (id);
\item разделители (delimiter, опционально);
\item неизвестные символы (unknown).
\end{enumerate}

\subsection{Форма Бэкуса–Наура (BNF)}

\begin{table}[H]
\centering
\renewcommand{\arraystretch}{1.35}
\begin{tabular}{ll}
\toprule
{\hfont\color{kviolet}Обычная запись} & {\hfont\color{kviolet}BNF}\\
\midrule
$\rul$ & $::=$\\
$S$ & $\langle S\rangle$\\
$S\rul A\mid AB$ & $\langle S\rangle::=A[B]$\\
$S\rul\alpha B',\ B'\rul\beta B'\mid\varepsilon$ & $\langle S\rangle::=\alpha\{\beta\}$\\
\bottomrule
\end{tabular}
\caption{Обычная запись правил и BNF}
\end{table}

Квадратные скобки — \term{необязательная часть}, фигурные — \term{повторение} (ноль и более раз). Обычно $\varepsilon$ в BNF не используют.

\section{Синтаксический анализ}

\subsection{Пример и классификация алгоритмов}

\begin{lstlisting}
if (a > b)
    a = a - b;
\end{lstlisting}
\[
\Downarrow\ \text{лексический анализатор}
\]
\begin{center}
\tok{if}\ \tok{(}\ \tok{a}\ \tok{>}\ \tok{b}\ \tok{)}\ \tok{a}\ \tok{=}\ \tok{a}\ \tok{-}\ \tok{b}\ \tok{;}
\end{center}
\[
\Downarrow\ \text{синтаксический анализатор}
\]

\begin{figure}[H]
\centering
\begin{forest}
  for tree={draw=kviolet,fill=kviolet!7,line width=0.9pt,rounded corners=2pt,font=\ttfamily,
    minimum height=6mm,minimum width=8mm,inner sep=2.5pt,l sep=6mm,s sep=7mm,edge={line width=0.8pt}},
  [if,draw=korange,fill=korange!12
    [>\vphantom{p} [a] [b]]
    [{=} [a] [{-}\vphantom{p} [a] [b]]]
    [nullptr,draw=kgray,fill=white,dashed]]
\end{forest}
\caption{Синтаксическое дерево: условие, тело и пустая ветка \texttt{else}}
\end{figure}

\begin{minipage}[t]{0.48\linewidth}
\textbf{Универсальные алгоритмы:}
\begin{enumerate}
\item КЯК (Кока–Янгера–Касами): $O(n^3)$, $n$ — длина ввода;
\item алгоритм Эрли: $O(n^3)$.
\end{enumerate}
\end{minipage}\hfill
\begin{minipage}[t]{0.48\linewidth}
\textbf{Неуниверсальные алгоритмы:}
\begin{enumerate}
\item $LL$-анализ (рекурсивный спуск);
\item $LR$-анализ.
\end{enumerate}
\end{minipage}

\subsection{Рекурсивный спуск}

\textbf{Пример.}
\[
\begin{array}{l}
\langle S\rangle::=\langle E\rangle\,\$\\
\langle E\rangle::=\langle T\rangle\ \{+\ \langle T\rangle\}\\
\langle T\rangle::=\langle F\rangle\ \{*\ \langle F\rangle\}\\
\langle F\rangle::=\langle Num\rangle\mid(\langle E\rangle)
\end{array}
\]
Разбираем строку: $5+7*(4+3*5)\$$.

\begin{figure}[H]
\centering
\begin{forest}
  for tree={font=\small,inner sep=1.5pt,l sep=3.5mm,s sep=2.5mm,edge={line width=0.7pt,kgray}},
  where n children=0{tier=word,text=korange,font=\ttfamily\bfseries}{text=kviolet,font=\small\bfseries},
  [S
    [E
      [T [F [5]]]
      [+]
      [T
        [F [7]]
        [*]
        [F
          [(]
          [E
            [T [F [4]]]
            [+]
            [T [F [3]] [*] [F [5]]]]
          [)]]]]
    [\$]]
\end{forest}
\hspace{1.2cm}
\begin{forest}
  for tree={circle,draw=korange,fill=korange!10,line width=0.9pt,minimum size=6.5mm,inner sep=0pt,font=\small,
    l sep=4mm,s sep=4mm,edge={line width=0.8pt}},
  where n children=0{draw=kviolet,fill=kviolet!7}{},
  [$+$ [$5$] [$*$ [$7$] [$+$ [$4$] [$*$ [$3$] [$5$]]]]]
\end{forest}
\caption{Дерево разбора строки $5+7*(4+3*5)\$$ (слева) и получаемое синтаксическое дерево (справа)}
\end{figure}

\textbf{Как писать.} $t$ — текущий токен, @gt()@ — взять следующий токен.

\textbf{Проверка:}
\begin{lstlisting}
void S() {
    E();
    if (t != '$') {
        throw false;
    }
    gt();
}

void E() {
    T();
    while (t == '+') {
        gt();
        T();
    }
}

void T() {
    F();
    while (t == '*') {
        gt();
        F();
    }
}

void F() {
    if (t == '(') {
        gt();
        E();
        if (t != ')') {
            throw false;
        }
        gt();
        return;
    }
    if (t.type == "num") {
        gt();
        return;
    }
    throw false;
}
\end{lstlisting}

\textbf{Построение:}
\begin{lstlisting}
Node* E() {
    Node* left = T();
    while (t == '+') {
        gt();
        Node* tmp = new Node('+');
        tmp->left = left;
        tmp->right = T();
        left = tmp;
    }
    return left;
}
\end{lstlisting}

\textbf{Для присваивания (правоассоциативный оператор):}
\[
\langle A\rangle::=\langle I\rangle\mid\langle I\rangle=\langle A\rangle
\]
\begin{lstlisting}
Node* A() {
    Node* left = I();
    if (t == '=') {
        gt();
        Node* tmp = new Node('=');
        tmp->left = left;
        tmp->right = A();
        left = tmp;
    }
    return left;
}
\end{lstlisting}

\subsection{Применимость рекурсивного спуска}

\begin{opr}
$\operatorname{First}(\alpha)$ — множество терминалов, с которых может начинаться $\alpha$, $\alpha\in(T\cup N)^*$.
\end{opr}

\textbf{Пример.}
\[
\begin{array}{l}
S\rul aA\mid BS\\ A\rul bB\mid cB\mid\varepsilon\\ B\rul aA\mid e
\end{array}
\qquad
\operatorname{First}(ABb)=\{b,c,a,e\},\quad\operatorname{First}(\varepsilon)=\varnothing.
\]

\begin{opr}
$\operatorname{Follow}(A)$ — множество терминалов, которые могут следовать за $A$, $A\in N$.
\end{opr}

\textbf{Пример.}
\[
\begin{array}{l}
S\rul AaS\mid Ab\\ A\rul aA\mid\varepsilon
\end{array}
\qquad
\operatorname{Follow}(A)=\{a,b\}.
\]

\begin{kvb}{$\operatorname{First}$ для нетерминалов $\operatorname{First}(A)$}
\begin{enumerate}
\item $A\rul a\alpha\Rightarrow$ добавляем $a$;
\item $A\rul B\alpha\Rightarrow$ добавляем $\operatorname{First}(B)$. Если $B\Rightarrow^*\varepsilon$, добавляем $\operatorname{First}(\alpha)$.
\end{enumerate}
\end{kvb}

\begin{kvb}{$\operatorname{First}$ для цепочек $\operatorname{First}(\alpha)$}
\begin{enumerate}
\item $\alpha=a\beta\Rightarrow\{a\}$;
\item $\alpha=A\beta\Rightarrow\operatorname{First}(A)$. Если $A\Rightarrow^*\varepsilon$, то $\operatorname{First}(A)\cup\operatorname{First}(\beta)$.
\end{enumerate}
\end{kvb}

\begin{kvb}{$\operatorname{Follow}(A)$}
\begin{enumerate}
\item $B\rul\xi A\alpha\Rightarrow$ добавляем $\operatorname{First}(\alpha)$;
\item $B\rul\xi A\Rightarrow$ добавляем $\operatorname{Follow}(B)$;
\item $B\rul\xi A\alpha$ и $\alpha\Rightarrow^*\varepsilon\Rightarrow$ добавляем $\operatorname{Follow}(B)$.
\end{enumerate}
\end{kvb}

\textbf{Пример.}
\[
\begin{array}{l}
S\rul aA\mid BDc\\ A\rul BAa\mid aB\mid b\\ B\rul\varepsilon\\ D\rul B\mid b
\end{array}
\qquad
\begin{array}{ll}
\operatorname{First}(B)=\varnothing, & \operatorname{Follow}(S)=\varnothing,\\
\operatorname{First}(D)=\{b\}, & \operatorname{Follow}(A)=\{a\},\\
\operatorname{First}(S)=\{a,b,c\}, & \operatorname{Follow}(B)=\{a,b,c\},\\
\operatorname{First}(A)=\{a,b\}, & \operatorname{Follow}(D)=\{c\}.
\end{array}
\]

\begin{kob}{Критерии применимости рекурсивного спуска}
Раскрываем $A$: $\alpha A\beta$ (нужно раскрытие единственным образом).
\begin{enumerate}
\item $A\rul\alpha\mid\beta$: \quad $\operatorname{First}(\alpha)\cap\operatorname{First}(\beta)=\varnothing$;
\item $A\Rightarrow^*\varepsilon$ и $A\rul\alpha_1\mid\dots\mid\alpha_n$: \quad $\exists!\,i:\ \alpha_i\Rightarrow^*\varepsilon$;
\item $A\Rightarrow^*\varepsilon$: \quad $\operatorname{First}(A)\cap\operatorname{Follow}(A)=\varnothing$.
\end{enumerate}
\end{kob}

\end{document}
